\subsection{Additional details regarding quark-sector modifications}
\label{ssec:theorydetails}
%
The minimal SME is grounded in the 
generic principles of effective field theory. 
It retains all of the SM properties, including
the $SU(3)\times SU(2)\times U(1)$ gauge structure
$SU(2)\times U(1)$ breaking, renormalizability, 
microcausality, spin-statistics, 
canonical quantization, energy-momentum conservation, and observer
Lorentz invariance. Its atypical properties are violations 
of particle Lorentz invariance and 
possibly CPT invariance
\footnote{Roughly half of the SME terms also violate CPT invariance.}.
The nonminimal SME relaxes the condition of renormalizability. 
As an effective field theory, the SME Lagrangian 
necessarily includes the SM Lagrangian, 
see Eq.~\eqref{SMEaction}. The additional Lorentz- and CPT-violating
terms are treated as perturbations, which is experimentally justified.
The SME is as a conservative extension of SM and
can be thought of as a 
``test framework" for Lorentz and CPT violation.

Our current focus is capturing the dominant 
signatures in the quark sector of the minimal SME.
To begin, consider a single-fermion field $\psi$.
Gauge interactions of $\psi$ include coupling constants.
As we work in perturbation theory, 
any Lorentz-violating effects tied to
coupling constants are expected to be 
suppressed relative to modified free-propagation 
effects, since the latter are absent of coupling 
constants. Loop contributions involve additional
powers of coupling constants and are further
suppressed. In this setting, the dominant
Lorentz- and CPT-violating effects 
on $\psi$ involve a generalized Dirac equation 
with several $d=3$ and $4$ operators~\cite{Colladay:2001wk}:
%
\begin{align}
\label{QEDSME}
\mathcal{L} = \tfrac{1}{2}\bar\psi i\Gamma^\mu \overset{\text{\tiny$\leftrightarrow$}}\partial_\mu\psi - \bar\psi M \psi,
\end{align}
%
where
%
\begin{align}
\label{gammaM}
&\Gamma^\mu = \gamma^\mu + c^{\nu\mu}\gamma_\nu + d^{\nu\mu}\gamma_5\gamma_\nu + e^\mu + if^\mu\gamma_5 + \tfrac{1}{2}g^{\alpha\beta\mu}\sigma_{\alpha\beta}, \\
&M = a_\mu\gamma^\mu + b_\mu \gamma_5\gamma^\mu + \tfrac{1}{2}H_{\mu\nu}\sigma^{\mu\nu}. 
\end{align}
%
The $-\bar\psi M\psi$ terms change sign under a CPT transformation 
(they are CPT odd). The operators are mass dimension $d=$
so the associated SME coefficients $a_\mu$, $b_\mu \gamma_5\gamma^\mu$
and $H^{\mu\nu}$ have units of GeV$^{-1}$. As such, 
in a dimensionless observable these terms scale like 
(SME coefficient)/energy and are dominant in low-energy processes.
Effects associated to these coefficients 
are expected to be small compared to
those contained in $\Gamma^\mu$ and
can be ignored here.
Within $\Gamma^\mu$, the coefficients
$e^\mu, f^\mu, g^{\alpha\beta\mu}$ are inconsistent
with the SM gauge structure and renormalizability~\cite{ck98}, 
and can also be ignored here. This
leaves the dimensionless and 
CPT-even $c$- and $d$-type discussed in the main text.
As a brief aside, though these coefficients are CPT even,
they introduce unconventional sources of other 
discrete-symmetry violations as shown in Tab.~\ref{summarycoeffs}.
%
\begin{table}[ht]
\caption{Coefficient properties under discrete symmetries
where $(-1)^\mu  = 1$ for $\mu = 0$ and $(-1)^\mu \equiv -1$ for $\mu = 1,2,3$.}
\centering 
\begin{tabular}{c c c c c c c c} % centered columns (4 columns)
\hline\hline
Coefficient & C & P & T & CP & CT & PT & CPT \\ [0.5ex] % inserts table
%heading
\hline % inserts single horizontal line
$c_{00}, c_{jk}$ & $+$ & $+$ & $+$ & $+$ & $+$ & $+$ & $+$\\
$c_{0j}, c_{j0}$ & $+$ & $-$ & $-$ & $-$ & $-$ & $+$ & $+$\\
$d_{00}, d_{jk}$ & $-$ & $-$ & $+$ & $+$ & $-$ & $-$ & $+$\\
$d_{0j}, d_{j0}$ & $-$ & $+$ & $-$ & $-$ & $+$ & $-$ & $+$\\[1ex] 
\hline\hline %inserts single line
\end{tabular}
\label{summarycoeffs}
\end{table}
%

Considering now $\psi$ as a quark field
in terms of $SU(2)$ doublets
and singlets, 
the $SU(2)\times U(1)$ invariant generalization
can be written as
%
\begin{align}
\label{model1}
\mathcal{L} &= \tfrac{1}{2}i c_{Q_i}^{\mu\nu}\overline{Q}_i\gamma_\mu\overset{\text{\tiny$\leftrightarrow$}}{D_{\nu}}Q_i  + \tfrac{1}{2}i c_{U_i}^{\mu\nu}\overline{U}_i\gamma_\mu\overset{\text{\tiny$\leftrightarrow$}}{D_{\nu}}U_i 
 + \tfrac{1}{2}i c_{D_i}^{\mu\nu}\overline{D}_i\gamma_\mu\overset{\text{\tiny$\leftrightarrow$}}{D_{\nu}}D_i \;,
\end{align}
%
where $D_\nu$ is the conventional 
gauge-covariant derivative. 
The fields have the usual definitions
%
\begin{align}
Q_i = \begin{pmatrix} u_i \\ d_i \end{pmatrix}_L, \quad U_i = \left(u_i\right)_R, \quad D_i = \left(d_i\right)_R.
\end{align}
%
The index $i = 1, 2, 3$ denotes flavors $u_i = (u, c, t)$ and 
$d_i = (d, s, b)$ and the SME coefficients are given 
by $c_{Q_i}^{\mu\nu}$, $c_{U_i}^{\mu\nu}$, and $c_{D_i}^{\mu\nu}$. 

Field redefinitions imply the antisymmetric parts of the coefficients 
are unobservable at first order in Lorentz 
violation~\cite{Colladay:1998fq, Fittante:2012ua, Colladay:2002eh}. 
The trace contributions are also unphysical because they can be 
absorbed into a rescaling of the fields. 
They are set to zero without loss of generality. 
This leaves nine observable components per coefficient type
$Q_i, U_i, D_i$. 
It is customary and useful to define 
combinations of these coefficients:
%
\begin{align}
\label{eq:quarkcoeffsapp}
c_{u_i}^{\mu\nu} &= (c_{Q_i}^{\mu\nu} + c_{U_i}^{\mu\nu})/2, \quad c_{d_i}^{\mu\nu} = (c_{Q_i}^{\mu\nu} + c_{D_i}^{\mu\nu})/2  \;, \nonumber\\
d_{u_i}^{\mu\nu} &= (c_{Q_i}^{\mu\nu} - c_{U_i}^{\mu\nu})/2, \quad d_{d_i}^{\mu\nu} = (c_{Q_i}^{\mu\nu} - c_{D_i}^{\mu\nu})/2  \;.
\end{align}
%
Since only three types of coefficients are initially present, 
the constraint 
$c_{u_i}^{\mu\nu} -c_{d_i}^{\mu\nu} 
= d_{d_i}^{\mu\nu} - d_{u_i}^{\mu\nu}$ eliminates
a redundancy in Eq.~\eqref{eq:quarkcoeffsapp}. 
Expressing Eq.~\eqref{model1} in terms of the $c$- and
$d$-type combinations introduced makes the comparison 
with the initial starting point Eq.~\eqref{QEDSME} more direct. 
%
\subsection{Comments on expected coefficient sensitivities}
%
Why are colliders suitable 
for probing the $c$- and $d$-type coefficients?
As discussed in Sec.~\ref{ssec:Context}, 
directly probing the properties of quarks
is often obscured by the strong interaction. 
However, the Drell-Yan process is one of a few 
where QCD factorization has been proven 
to all orders in the strong coupling constant. 
This enables analytical predictions for
the scattering cross sections in terms
of partonic matrix elements convoluted
with parton distribution functions to be obtained. 

An alternative path is to invoke 
effective theories like chiral perturbation
theory ($\chi$PT) to relate
parton-level coefficients to effective meson 
and baryon coefficients~\cite{Altschul:2019beo}.
One could then use existing constraints 
on the effective proton and neutron
coefficients $c_{p,n}^{\mu\nu}$ from, e.g.
atomic-clock frequency comparisons to constrain
the linear combinations of (valence) quark coefficients
underlying the hadron coefficients.
For example, in the absence of gluonic modifications,
one obtains an effective $c$-type
proton coefficient \cite{Kamand:2016xhv}
\begin{align}
c_p^{\mu\nu} = \left[\frac{1}{2}\alpha^{(1)} + \alpha^{(2)}\right]
(c_{uL}^{\mu\nu}+ c_{uR}^{\mu\nu}) 
+ \left[-\frac{1}{2}\alpha^{(1)} + \alpha^{(2)}\right]
(c_{dL}^{\mu\nu}+c_{dR}^{\mu\nu}),
\end{align}
where $\alpha^{(1)}, \alpha^{(2)}$ are low-energy constants.
These coefficients are currently unknown, but could 
in principle be calculated using lattice QCD.
Such route is not presently available and 
constraints on the quark-level parameters 
from $\chi$PT are technically unavailable. 
That said, by invoking considerations of naturalness, 
one could estimate that the quark constraints
should be of a comparable order of magnitude to the hadron ones,
unless some dramatic cancellation occurs. 

On the extremely high-energy scale, and as 
mentioned in Sec.~\ref{ssec:prevstudies}, 
astrophysical measurements have been suggested
as potential probes of the $c$-type light-quark ($u,d$)
coefficients. For the top quark, 
some constraints from threshold analyses 
imply the constraint
$|c_t^{\mu\nu}| < 10^{-4}$, 
surpassing the those from $t\bar t$
production by $\sim$2-3 orders of magnitude. 
Projections from upgrades at the LHC 
and the FCC claim to be able to surpass
those from threshold analyses
by 1-2 orders of magnitude~\cite{Carle:2019ouy}.
Phenomenological constraints from radiative corrections
have also been placed. Only the $c$-type coefficients
have the correct tensor structure 
and symmetry properties
for contributing to the one-loop 
photon vacuum polarization, leading to 
constraints on the purely timelike component
$|c_t^{00}| \simeq 10^{-7}$~\cite{Satunin:2017wmk}. 

In contrast, alternative probes of 
the $d$-type quark coefficients do not exist.
The central complication is 
that these coefficients affect the quarks' polarization structure.
This leads to significantly fewer observables 
and straightforward experimental probes.
Virtual gluonic effects inside
a bound hadron imply a given parton 
does not have a well-defined spin.
For a spin-zero meson, the net spin from 
valence quarks must 
average to zero, rendering the $d$-type effects
are unobservable in many processes~\cite{Altschul:2020wkw}.
Reiterating some comments made in 
Sec.~\ref{ssec:LIVDY},
the theory STA members 
have shown~\cite{ls18} that
unpolarized hadronic processes
mediated by a photon or gluon
are completely insensitive to $d$-type coefficients.
This apparent obstruction could be surmounted
with polarized beamlines, but no existing
collider has this capability.
The key insight is to consider processes
mediated instead by massive gauge bosons
with distinct of left- and right-handed
quark couplings~\cite{lssv21}. 
This is the only known route to 
probe the $d$-type quark coefficients.
The Drell-Yan process is a perfect test process
because of the wealth of existing data 
at invariant masses $Q\sim m_Z$, where
the $d$-type coefficients display the 
largest cross-section enhancement.
Deep inelastic scattering with $t$-channel $Z$ exchange
is another possibility, but is greatly suppressed.

In summary, with regard to the
$c$-type coefficients, the 
expected constraints~\cite{ls18,klsv19}
will not be competitive with existing 
first-generation ($u,d$) results, however,
new linear coefficient combinations will be tested.
Second-generation ($s,c$)
constraints will exceed, 
by a few orders of magnitude~\cite{lssv21},
the few existing constraints~\cite{Karpikov:2016qvq}.
The $d$-type coefficients are much more
motivated and suitable for our analysis.
